% =====================================================================
%  Section 1: Introduction
% =====================================================================
\section{Introduction}
\label{sec:intro}

\paragraph{Data selection in instruction tuning.}
Instruction tuning (IT) is the foundational stage of LLM
post-training: contemporary open releases~\cite{dubey2024llama3,
lambert2024tulu3,qwen2025} begin with supervised fine-tuning on
instruction--response pairs. As IT corpora have grown from thousands
of pairs to hundreds of thousands, an empirical near-consensus has
emerged: at fixed compute, the quality and composition of the IT
subset dominate its raw scale~\cite{zhou2023lima,zhang2025jair,
qin2024datatsunami}. The central lever of IT is therefore
\emph{selection}---choosing which subset of a large candidate pool
to train on.

\paragraph{Selection methods.}
Selection criteria have evolved along three model-aware
axes---surface-level quality scoring, gradient influence on a
target validation set, and structural signals from hidden-state
geometry (Sec.~\ref{sec:related}). A more recent line is
\emph{dynamic data selection}: re-estimating data utility during
training as the model state changes~\cite{kung2023activeit,
sun2025dots}. This line shows that useful examples are not fixed
properties of the data alone; their value can depend on the current
model, training stage, and optimisation regime. In this paper, we
focus on the supervised instruction-tuning setting, where candidates
are instruction--response pairs and no rollout group, verifier
reward, or target validation set is assumed.

\paragraph{Two gaps in dynamic selection.}
Despite this progress, two gaps remain. First, structural
hidden-state selectors and dynamic selectors have largely evolved
separately. The former extract capability directions from model
representations, but usually fix them before training using seed or
target examples. The latter adapt during training, but typically
optimise scalar utilities such as difficulty, uncertainty, loss
change, or reward proxies without directly using the model's
evolving representation geometry. As a result, dynamic selection can
track the training state while still missing the structural direction
along which the model is being updated.

Second, the scalar utilities used in dynamic selection can be
unstable in multi-task instruction pools. In the RL composite-reward
selector we use as a base, difficulty and uncertainty are combined
through a variance-ratio weight. In multi-task pools, this weight can
shift with task composition because it conflates within-task and
between-task variance; under a naive per-micro-batch implementation,
it can even lose the difficulty signal when the micro-batch size is
one. These sensitivities motivate a dynamic selection criterion that
is not only adaptive, but also representation-aware and stable.

\paragraph{Our approach: Trajectory-Anchored Dynamic Selection.}
We propose \textsc{TADS} (Trajectory-Anchored Dynamic Selection), a
model-intrinsic framework that augments a dynamic selection pipeline
with a structural signal extracted from the model's own hidden-state
trajectory. At each refresh step, \textsc{TADS} computes
contextualisation deltas
\[
\Delta h_l(x;\theta_{t-1})
=
h_l^{\mathrm{last}}(x;\theta_{t-1})
-
h_l^{\mathrm{first}}(x;\theta_{t-1})
\]
on a probe subset of the candidate pool and takes the top principal
direction as a per-layer trajectory anchor $v_l^{(t)}$. This anchor
requires no in-domain seed examples, no validation set, no labels,
and no auxiliary scorer. It is extracted from the same model being
fine-tuned, using the candidate pool itself.

The anchor enters the selection rule through a single multiplicative
factor. The underlying advantage-estimator pipeline provides two
time-dependent signals for each candidate at refresh step $t$: a
composite reward $R_i^{(t)}$ and a learned advantage estimate
$a_i^{(t)} \in (0,1)$ (App.~\ref{app:setup}). We use their product
as a \emph{reward--advantage diagnostic consensus} score: a
candidate receives a high base score only when both the scalar
reward signal and the advantage estimate support its selection.
The advantage estimate $a_i^{(t)}$ is the \emph{output} of a learned
estimator; the training procedure used to fit it
(App.~\ref{app:setup}) is incidental to the selection rule
itself, which consumes $a_i$ only as a bounded scalar value.
\textsc{TADS} then applies the trajectory anchor as
\[
s_i^{(t)}
=
R_i^{(t)} a_i^{(t)}
\left(
1+\lambda\,\widetilde{\mathrm{align}}_i^{(t)}
\right),
\]
where $\widetilde{\mathrm{align}}_i^{(t)}\in[0,1]$ measures how well
the candidate's sequence-mean hidden state aligns with the current
trajectory anchor, averaged across layers. Setting $\lambda=0$
removes the anchor and recovers the fixed reward--advantage consensus
base ranking $R_i^{(t)}a_i^{(t)}$; varying $\lambda$ therefore
isolates the contribution of the trajectory anchor on top of the
same base rule.

\paragraph{Stability of the trajectory anchor.}
A key question is whether the anchor is a meaningful structural
object or merely probe noise. We answer this with a stability
guarantee. Under a positive eigenvalue gap and bounded covariance
change, the sign-calibrated trajectory anchor satisfies
\[
\left\|
v_l^{(t+1)}-v_l^{(t)}
\right\|
\le
\frac{2C_\Sigma}{\gamma_{\min}}\,\eta_t .
\]
Thus the anchor evolves smoothly at the scale of the learning-rate
schedule. The guarantee concerns the trajectory anchor itself and
therefore the structural component introduced by \textsc{TADS}; it
does not depend on interpreting the reward--advantage consensus score as
a policy-gradient estimator. This separation is important: the base
ranking supplies a fixed diagnostic selection rule, while the theorem
establishes stability for the additional trajectory-derived signal.

\paragraph{Contributions.}
This paper makes four contributions. First, we revisit dynamic
instruction-tuning data selection and identify a practical
instability in composite-reward selection: the variance-ratio weight
can drift with task mix and can collapse under naive micro-batch
statistics. Second, we introduce \textsc{TADS}, a trajectory-anchored
dynamic selector that converts hidden-state contextualisation deltas
into per-layer capability directions and uses them as a
$\lambda$-controlled multiplicative anchor on top of a fixed
reward--advantage diagnostic consensus base. Third, we establish a
stability guarantee for the trajectory anchor, showing that its
epoch-to-epoch drift is controlled by the learning-rate schedule
under a positive eigenvalue gap; the guarantee concerns the
structural anchor rather than any particular scalar reward or advantage
ranking rule. Finally, we empirically evaluate \textsc{TADS} across
model scales, instruction-tuning datasets, and benchmark families,
showing improvements over strong static, seed-based, and dynamic
baselines with less than $0.5\%$ wall-clock overhead, while
$\lambda$ ablations isolate the contribution of the trajectory
anchor.
